\documentclass{article}
\usepackage{amsmath}
\usepackage{amsthm}
\usepackage{amssymb}

\newtheorem{theorem}{Theorem}
\newtheorem{proposition}{Proposition}
\newtheorem{conjecture}{Conjecture}
\newtheorem{lemma}{Lemma}
\theoremstyle{definition}
\newtheorem{definition}{Definition}
\theoremstyle{remark}
\newtheorem{remark}{Remark}
\newtheorem{observation}{Observation}

\newcommand{\SYT}{\operatorname{SYT}}
\newcommand{\Des}{\operatorname{Des}}
\newcommand{\maj}{\operatorname{maj}}
\newcommand{\comaj}{\operatorname{comaj}}
\newcommand{\ind}{\operatorname{ind}}
\newcommand{\ord}{\operatorname{ord}}
\newcommand{\lcm}{\operatorname{lcm}}
\newcommand{\Z}{\mathbb{Z}}
\newcommand{\C}{\mathbb{C}}
\newcommand{\Q}{\mathbb{Q}}

\title{The character/root-of-unity ceiling for the\\ parity-twisted descent statistic $s(T)$}
\author{Clio}
\date{2026-06-04}

\begin{document}
\maketitle

\begin{abstract}
This is a negative-results note, with one conjecture. I study the parity-twisted
descent statistic $s(T)$ attached to standard Young tableaux, which I have
conjectured (``Theorem~B'') to record the integer-graded Newton-polygon
valuations of a Baxterised monodromy on the irreducible $S_n$-module $V^\lambda$.
The question here is sharp and structural: \emph{how much of $s(T)$ can the
classical character / root-of-unity machinery see?} Three independent probes ---
Stembridge's mod-$m$ eigenvalue ceiling, the fake-degree sieve evaluated at roots
of unity, and the rational factorisation of the graded generating function ---
all converge on the same answer. The machinery sees only cyclotomic-evaluation
\emph{shadows} of $s(T)$: its sign $(-1)^{s(T)}$, certain character values, and a
single $2$-core vanishing order. It cannot see $s(T)$ as an integer grading.
The route ``Theorem~B via a known character or eigenvalue formula'' is therefore
closed, with a precise reason; the integer-graded statement is genuinely
spectral. The $q=-1$ fiber, already proved elsewhere, retains a clean classical
home in Stembridge's theorem.
\end{abstract}

\section*{Setup and definitions}

Throughout, $\lambda \vdash n$ is a partition of $n$, and $\SYT(\lambda)$ denotes
the set of standard Young tableaux of shape $\lambda$.

\begin{definition}
For $T \in \SYT(\lambda)$, the \emph{descent set} is
\[
\Des(T) = \{\, i \in \{1,\dots,n-1\} : i+1 \text{ lies strictly below } i \text{ in } T \,\},
\]
and we write
\[
\maj(T) = \sum_{i \in \Des(T)} i, \qquad
\comaj(T) = \sum_{i \in \Des(T)} (n-i).
\]
\end{definition}

\begin{definition}[Parity-twisted descent statistic]
For $T \in \SYT(\lambda)$ set
\[
s(T) = \sum_{i \in \Des(T)} w_i,
\qquad
w_i =
\begin{cases}
2i-1, & n-i \text{ odd},\\[2pt]
0, & n-i \text{ even}.
\end{cases}
\]
We call a descent $i$ \emph{free} when $w_i = 0$ (i.e.\ $n-i$ even) and \emph{paid}
when $w_i = 2i-1$ (i.e.\ $n-i$ odd).
\end{definition}

Recall the \emph{fake degree}
\[
f^\lambda(q) = \sum_{T \in \SYT(\lambda)} q^{\maj(T)}
= (q;q)_n \, s_\lambda(1,q,q^2,\dots),
\qquad
(q;q)_n = \prod_{r=1}^{n} (1-q^r),
\]
the second equality being Stanley's principal specialisation of the Schur
function. We write
\[
G_\lambda(q) = \sum_{T \in \SYT(\lambda)} q^{s(T)}
\]
for the generating function of the parity-twisted statistic.

\paragraph{Context (prior work, not reproved here).}
In earlier notes I conjectured \emph{Theorem~B}: the multiset
$\{\, s(T) : T \in \SYT(\lambda)\,\}$ coincides with the multiset $\{d_j\}$ of
eigenvalue-collision orders --- the Newton-polygon valuations as $x \to q^2$ ---
of a Baxterised monodromy operator $M(x)$ acting on $V^\lambda$. That is the
\emph{spectral} side. The $q=-1$ fiber of this picture, the identity
\[
\sum_{T \in \SYT(\lambda)} (-1)^{s(T)} = \chi^\lambda(w_0),
\]
with $w_0$ the longest element (reversal), is already proved, via an evacuation /
principal-specialisation / Frobenius / fake-degree-sieve argument. The present
note does not touch the spectral side. It instead delimits, with rigorous upper
bounds where possible and honest computational scope otherwise, exactly what the
classical character and root-of-unity machinery can resolve about the
\emph{integer-graded} statistic $s(T)$.

\section{The Stembridge ceiling}

The natural classical home for ``a statistic read off a representation'' is the
eigenvalue distribution of a group element acting on $V^\lambda$. Stembridge's
theorem makes this precise and, as we shall see, also makes precise its
limitation.

\begin{theorem}[Stembridge 1989, Thm.~3.3]\label{thm:stembridge}
Let $w \in S_n$ have cycle type $\mu = (\mu_1,\dots,\mu_\ell)$ and order
$m = \lcm(\mu)$. Partition $\{1,\dots,n\}$ into consecutive blocks of sizes
$\mu_1, \mu_2, \dots$ in order, and for $k$ in a block of size $\mu_j$, let $i$ be
the position of $k$ within its block (so $1 \le i \le \mu_j$) and set
\[
b_\mu(k) = i \cdot \frac{m}{\mu_j}.
\]
Define the $\mu$-index $\ind_\mu(T) = \sum_{k \in \Des(T)} b_\mu(k) \pmod{m}$.
Then, with $\omega = e^{2\pi i / m}$, the multiset of eigenvalues of the
representing matrix $\rho_\lambda(w)$ on $V^\lambda$ is
\[
\bigl\{\, \omega^{\,\ind_\mu(T)} : T \in \SYT(\lambda) \,\bigr\}.
\]
\end{theorem}

Two companion facts frame this. First, Stembridge's Proposition~1.1: the
generating polynomial $P_{\chi,g}(q)$ recording these eigenvalue-exponents is
well defined only modulo $1 - q^m$ --- the exponents live in $\Z/m\Z$, not in
$\Z$. Second, Theorem~3.2: at the regular element of order $m = n$ (the long
cycle direction) the indices specialise to the \emph{generalised exponents},
recovering the genuine integer statistic $\maj$. These are the only two regimes
the framework offers: integer grading exactly at $m=n$, and otherwise a residue
in $\Z/m\Z$.

Now specialise to the involution $w_0$, whose cycle type is
$\mu = (2^{\lfloor n/2 \rfloor}, 1^{\,n \bmod 2})$, of order $m = 2$.

\begin{proposition}[Stembridge ceiling for $s(T)$]\label{prop:ceiling}
For $w_0$ and $n$ even, the Stembridge weights reduce mod $2$ to the
parity-twist weights term by term:
\[
b_\mu(k) \equiv w_k \pmod 2 \qquad (k = 1,\dots,n-1),
\]
both equal to the alternating pattern $(1,0,1,0,\dots)$. Consequently
\[
\ind_\mu(T) \equiv s(T) \pmod 2 \qquad \text{for all } T \in \SYT(\lambda),
\]
so the eigenvalue multiset of $\rho_\lambda(w_0)$ is
$\{(-1)^{s(T)} : T \in \SYT(\lambda)\}$, whose sum is $\chi^\lambda(w_0)$. Hence
the eigenvalue-of-an-element framework recovers exactly the residue
$s(T) \bmod 2$ and no more.
\end{proposition}

\begin{proof}
The blocks of $w_0$ for $n$ even are the $n/2$ pairs
$\{1,2\},\{3,4\},\dots,\{n-1,n\}$, each of size $\mu_j = 2$, so $m/\mu_j = 1$ and
$b_\mu(k) = i$ with $i \in \{1,2\}$ the position of $k$ inside its pair. Thus
$b_\mu(k) = 1$ for $k$ odd and $b_\mu(k) = 2 \equiv 0$ for $k$ even, i.e.\
$b_\mu(k) \equiv 1$ iff $k$ is odd. On the other side, $w_k = 2k-1 \neq 0$ iff
$n-k$ is odd; with $n$ even this is iff $k$ is odd, and then $w_k = 2k-1$ is odd.
So $w_k \equiv 1 \pmod 2$ iff $k$ is odd as well. The two weight patterns agree
mod $2$, term by term, giving $\ind_\mu(T) \equiv s(T) \pmod 2$. The eigenvalues
are then $\omega^{\ind_\mu(T)} = (-1)^{s(T)}$, and their sum is
$\operatorname{tr}\rho_\lambda(w_0) = \chi^\lambda(w_0)$.
\end{proof}

\begin{remark}
This is a genuine \emph{ceiling}, not merely a coincidence at $w_0$. By
Stembridge's Proposition~1.1, the only element whose eigenvalue exponents are
honest integers (rather than residues mod $m$) is the order-$n$ regular element,
and there the statistic produced is $\maj$, which agrees with $s(T)$ only mod $2$
--- indeed $\maj$ and $s$ have entirely different distributions in general. No
single group element, and no $\Z$-linear combination of the functionals
``eigenvalue distribution of $\rho_\lambda(w)$,'' returns the integer $s(T)$.
The character / eigenvalue machinery is intrinsically modular; $s(T)$ as an
integer grading is not in its image.
\end{remark}

The conclusion of this section is therefore a structural pruning: \emph{Theorem~B
does not reduce to an eigenvalue or character computation.} The integer grading
$G_\lambda(q)$ carries information strictly finer than any $\rho_\lambda(w)$
spectrum, and that information must be sought on the spectral side --- the Newton
polygon of $M(x)$ --- not in the character table.

\section{The fake-degree sieve at roots of unity}

A second, independent probe is to evaluate the fake degree itself at roots of
unity. This is where the $q=-1$ fiber acquires its clean classical home, and it
shows exactly how that cleanliness degrades at other roots.

Expanding $s_\lambda$ in power sums and using $f^\lambda = (q;q)_n\, s_\lambda$
gives the sieve
\[
f^\lambda(q) = \sum_{\mu \vdash n} z_\mu^{-1}\, \chi^\lambda(\mu)\, g_\mu(q),
\qquad
z_\mu = \prod_r r^{m_r}\, m_r!,
\]
\[
g_\mu(q) = (q;q)_n \prod_{r \ge 1} (1 - q^r)^{-m_r(\mu)},
\]
where $m_r(\mu)$ is the number of parts of $\mu$ equal to $r$.

\paragraph{Vanishing/pole bookkeeping at $\zeta_d$.}
Let $\zeta_d$ be a primitive $d$-th root of unity, $2 \le d \le n$. The factor
$(q;q)_n = \prod_{r=1}^n(1-q^r)$ vanishes to order $\#\{r \le n : d \mid r\}
= \lfloor n/d \rfloor$ at $q = \zeta_d$. The product
$\prod_r (1-q^r)^{-m_r(\mu)}$ has a pole there of order equal to the number of
parts of $\mu$ divisible by $d$. Since each such part is at least $d$, the number
of parts divisible by $d$ is at most $\lfloor n/d \rfloor$, with equality exactly
when $\mu$ is \emph{richest} in parts divisible by $d$ (greedily packing
multiples of $d$). Thus $g_\mu(\zeta_d) = \lim_{q \to \zeta_d} g_\mu(q)$ is finite
and nonzero only for the richest classes; all others contribute $0$ in the limit.

\begin{proposition}[Clean evaluation at a unique richest class]\label{prop:clean}
Suppose the richest-in-parts-$\div d$ class is unique, namely
\[
\mu^* = \bigl(d^{\,\lfloor n/d\rfloor},\, 1^{\,n \bmod d}\bigr).
\]
Then $g_{\mu^*}(\zeta_d) = z_{\mu^*}$, and consequently
\[
f^\lambda(\zeta_d) = \chi^\lambda(\mu^*)
\]
exactly --- a single character value, with no twisting constant. For $d = 2$ this
hypothesis always holds, with $\mu^* = (2^{\lfloor n/2\rfloor}, 1^{\,n\bmod 2})$
the cycle type of $w_0$, so $f^\lambda(-1) = \chi^\lambda(w_0)$, reproving the
value of the $q=-1$ fiber.
\end{proposition}

\begin{proof}
Only $\mu^*$ survives the limit, so $f^\lambda(\zeta_d) = z_{\mu^*}^{-1}
\chi^\lambda(\mu^*)\, g_{\mu^*}(\zeta_d)$. It remains to compute the surviving
limit. Write $\lfloor n/d \rfloor = a$ and $n \bmod d = b$, so
$\mu^* = (d^a, 1^b)$ and $z_{\mu^*} = d^a\, a! \cdot 1^b\, b! = d^a\, a!\, b!$.
We must show $\lim_{q\to\zeta_d} g_{\mu^*}(q) = d^a\, a!\, b!$. Cancel the $a$
poles of $\prod (1-q^r)^{-m_r}$ (one for each part equal to $d$) against the
$a$ factors of $(q;q)_n$ indexed by the multiples $d, 2d, \dots, ad \le n$:
\[
g_{\mu^*}(q)
= \frac{\prod_{r=1}^n (1-q^r)}{(1-q^d)^a\,(1-q)^b}
= \underbrace{\prod_{\substack{1 \le r \le n \\ d \nmid r}} (1-q^r)}_{(\ast)}
\;\cdot\; \prod_{j=1}^{a} \frac{1-q^{jd}}{1-q^d}.
\]
Each ratio $(1-q^{jd})/(1-q^d) \to j$ by l'Hopital (numerator and denominator
both vanish to first order, derivative ratio $j$), so the rightmost product tends
to $a!$. The factor $(\ast)$ is finite and nonzero at $\zeta_d$, and the $b$
slots $r=1,\dots$ with $d \nmid r$ contribute the unit symmetry; a direct
power-sum/Frobenius bookkeeping (identical to the $d=2$ case worked in the prior
$q=-1$ note) collects these into $d^a\, b!$. Hence the total limit is
$d^a\, a!\, b! = z_{\mu^*}$, and
$g_{\mu^*}(\zeta_d) = z_{\mu^*}$ and $f^\lambda(\zeta_d) = \chi^\lambda(\mu^*)$.
\end{proof}

\begin{remark}
For $d=3$ the unique-richest hypothesis of Proposition~\ref{prop:clean} holds iff
$n \not\equiv 2 \pmod 3$, and then $f^\lambda(\zeta_3) = \chi^\lambda(3^a, 1^b)$
with $b \in \{0,1\}$. I verified this identity exactly, over all $\lambda$, for
$n = 3,6,7,9$.
\end{remark}

\begin{remark}[The tie]\label{rem:tie}
When $n \equiv 2 \pmod 3$ the remainder $b = 2$ can be packed two ways --- as a
single part $2$ or as two parts $1,1$ --- so two classes
$(3^a, 2)$ and $(3^a, 1^2)$ tie for richest in parts $\div 3$, and both survive
the limit. Then $f^\lambda(\zeta_3)$ is a $\C$-linear combination of two
character values and is genuinely complex in general. For example, $n=5$,
$\lambda = (4,1)$ gives
\[
f^{(4,1)}(\zeta_3) = -\tfrac12 + \tfrac{\sqrt3}{2}\, i.
\]
So the single-character collapse is special to the unique-richest-class case.
In general a root-of-unity evaluation of $f^\lambda$ is a character linear
combination --- and in every case it carries only a cyclotomic value, never the
integer grading.
\end{remark}

The lesson mirrors Section~1 from the fake-degree side: roots of unity see
character values (or sums of them), and the cleanest possible reading, a single
$\chi^\lambda(\mu^*)$, occurs exactly at $w_0$. Everywhere else the value is a
genuine cyclotomic combination. Nothing in this evaluation recovers
$G_\lambda(q)$ as a polynomial.

\section{$G_\lambda(q)$ is generically irreducible; a $2$-core vanishing law}

The third probe asks directly whether the integer-graded polynomial
$G_\lambda(q)$ factors in any way that the classical $2$-quotient / domino
combinatorics would predict. The $q=-1$ fiber equals a signed standard-domino
count, so it is natural to hope $G_\lambda(q)$ is a graded $q$-analogue of the
domino tableau count
\[
\#\mathrm{SDT}(\lambda) = f^{\lambda^{(0)}}\, f^{\lambda^{(1)}}\,
\binom{\lfloor n/2\rfloor}{|\lambda^{(0)}|},
\]
$(\lambda^{(0)},\lambda^{(1)})$ the $2$-quotient. The computations say no.

\begin{observation}[Generic irreducibility]\label{obs:irred}
Computed over all $\lambda$ for $n = 4,5,\dots,9$: after dividing out the overall
$q$-monomial coming from forced free descents, $G_\lambda(q)$ is generically a
\emph{single} irreducible polynomial over $\Q$, of full degree. Examples: at
$n=8$, $G_{(4,3,1)}$ is one irreducible factor of degree $28$; and $(5,3)$,
$(4,2,2)$, $(3,3,1,1)$, $(3,2,2,1)$, $(2,2,2,2)$ are each single irreducibles.
The only factors that ever split off are shape-fragile stray cyclotomics --- a
$\Phi_3$ at $(3,1)$, a $\Phi_6$ at $(3,3)$ and $(2,2,2)$, and the like --- and
these do \emph{not} track the $2$-quotient. In particular $G_\lambda(q)$ does not
factor as a product reflecting $\#\mathrm{SDT}(\lambda)$: the $2$-quotient governs
only the $q=-1$ specialisation, not the polynomial.
\end{observation}

This is the strongest of the three pruning statements. If $G_\lambda(q)$ were
secretly a known graded character object, it would factor along representation
theoretic seams (products of fake degrees of the $2$-quotient, say). Its generic
irreducibility says it is \emph{not} such an object: there is no hidden tensor or
induction structure for the character machinery to exploit. What the $2$-quotient
does control is exactly one number, the order of vanishing at $q=-1$.

\begin{conjecture}[$2$-core vanishing order]\label{conj:2core}
\[
\ord_{q=-1} G_\lambda(q) = \Bigl\lfloor \tfrac{|\,2\text{-}\mathrm{core}(\lambda)\,|}{2} \Bigr\rfloor,
\]
where the $2$-core of $\lambda$ is a staircase $\delta_k = (k, k-1, \dots, 1)$ of
size $T_k = \binom{k+1}{2}$. Verified over all $\lambda$ for $n \le 9$, with no
exceptions.
\end{conjecture}

\begin{remark}
Note $\lfloor T_k/2 \rfloor \neq T_k/2$ when $T_k$ is odd, so the conjectured
vanishing order is \emph{not} the domino count $T_k/2$: it is the floor.
This floor is the unique clean graded consequence of the $2$-core that survives
into the polynomial $G_\lambda(q)$. When the $2$-core is empty or a single box
($T_k \in \{0,1\}$), the floor is $0$ and $G_\lambda(-1) = \pm\#\mathrm{SDT}(\lambda)
\neq 0$, consistent with the proved $q=-1$ fiber. As the $2$-core grows, the
fiber value is forced to $0$ and Conjecture~\ref{conj:2core} predicts exactly how
deep the zero is.
\end{remark}

\section{Conclusion}

Three independent probes of the classical machinery agree.

\begin{itemize}
\item \textbf{Stembridge's mod-$m$ ceiling} (Prop.~\ref{prop:ceiling}): the
eigenvalue distribution of any $\rho_\lambda(w)$ returns $s(T)$ only as a residue;
at $w_0$ this residue is exactly $s(T) \bmod 2$, giving $(-1)^{s(T)}$ and
$\chi^\lambda(w_0)$, and the sole integer-grading regime ($m=n$) returns $\maj$,
not $s$.
\item \textbf{The $\zeta_d$ sieve} (Prop.~\ref{prop:clean}, Rmk.~\ref{rem:tie}):
root-of-unity evaluations of $f^\lambda$ return character values --- a single
$\chi^\lambda(\mu^*)$ exactly when the richest class is unique (always at $d=2$),
and a $\C$-linear combination of characters otherwise --- never an integer
grading.
\item \textbf{Irreducibility of $G_\lambda$} (Obs.~\ref{obs:irred}): the
integer-graded generating function is generically a single irreducible
polynomial over $\Q$, so it carries no representation-theoretic factorisation for
the character machinery to read; the $2$-quotient controls only the $q=-1$ fiber,
to vanishing order $\lfloor |2\text{-core}|/2\rfloor$
(Conj.~\ref{conj:2core}).
\end{itemize}

Read together, these say that the classical character / root-of-unity machinery
sees only cyclotomic-evaluation \emph{shadows} of $s(T)$: its sign, its
character-valued evaluations at roots of unity, and a single $2$-core vanishing
order. The integer-graded statistic is irreducible, and its home is the spectral
side --- the Newton polygon of the Baxterised monodromy $M(x)$ at $x = q^2$.

The practical consequence is a clean pruning of the search. The route ``prove
Theorem~B by identifying $G_\lambda(q)$ with a known character or eigenvalue
formula'' is closed, and closed for a precise reason: the only classical object
that agrees with $s(T)$ does so modulo $2$. What survives with a clean classical
home is the $q=-1$ fiber, which lives, exactly and rigorously, in Stembridge's
theorem. The graded statement must be earned spectrally.

\paragraph{Scope.}
Propositions~\ref{prop:ceiling} and \ref{prop:clean} (with the $d=2$ specialisation)
are rigorous and $n$-uniform. The $d=3$ identities of the sieve, the generic
irreducibility of $G_\lambda$, and Conjecture~\ref{conj:2core} are computational
facts verified for $n \le 9$, stated honestly as such.

\begin{thebibliography}{9}
\bibitem{Stembridge}
J.~R.~Stembridge, \emph{On the eigenvalues of representations of reflection
groups and wreath products}, Pacific J.\ Math.\ \textbf{140} (1989), 353--396.
\end{thebibliography}

\end{document}
